%\clearpage
\section{Funciones monótonas}

\begin{definicion}
Sea $f$ una función real definida en un subconjunto $A\subseteq\R$.
\begin{enumerate}
    \item $f$ es creciente (no decreciente) en $A$ si $x<y$ implica $f(x)\le f(y)$.
    \item $f$ es estrictamente creciente en $A$ si $x<y$ implica $f(x)<f(y)$.
    \item $f$ es decreciente (resp. estrictamente decreciente) si $x<y$ implica $f(x)\ge f(y)$ (resp. $f(x)>f(y)$).
    \item $f$ es monótona en $A$ si es creciente o decreciente.
\end{enumerate}
\end{definicion}

\begin{teorema}
Sea $f:(a,b)\to\R$ una función monótona en $(a,b)$ y $c\in(a,b)$.
Entonces existen $f(c^-)$ y $f(c^+)$.

Además, si $f$ es creciente, entonces
\[
f(c^-)=\sup\{f(x):a<x<c\}\le f(c)\le \inf\{f(x):c<x<b\}=f(c^+).
\]
Si $f$ es decreciente, se invierten las desigualdades:
\[
f(c^-)=\inf\{f(x):a<x<c\}\ge f(c)\ge \sup\{f(x):c<x<b\}=f(c^+).
\]
\end{teorema}

\begin{corolario}
Sea $f:(a,b)\to\R$ monótona. Entonces $f$ no tiene discontinuidades de segunda clase.
\end{corolario}

\begin{teorema}
Sea $f$ monótona en $(a,b)$. El conjunto de puntos de discontinuidad de $f$ en $(a,b)$ es a lo sumo numerable.
\end{teorema}

\begin{teorema}
Si $f$ es estrictamente creciente en $A\subseteq\R$, entonces $f^{-1}$ existe y es estrictamente creciente en $f(A)$.
\end{teorema}

\begin{teorema}
Sea $f$ estrictamente creciente y continua en un intervalo compacto $[a,b]$.
Entonces $f^{-1}$ es estrictamente creciente y continua en $[f(a),f(b)]$.
\end{teorema}
